\documentclass[11pt]{article}
\usepackage[margin=1in]{geometry}
\usepackage{amsmath,amssymb,amsthm}
\usepackage{booktabs}
\newtheorem{proposition}{Proposition}
\newtheorem{definition}{Definition}
\newtheorem{remark}{Remark}
\DeclareMathOperator{\tr}{tr}
\begin{document}

\section*{Step 1 Structural Statement: Mode A Holographic Import Audit}

\begin{definition}[Leading recognition source]
For Step 1, the leading named recognition source is the holographic
entanglement dictionary: the AdS/CFT dictionary together with the
Ryu--Takayanagi / Hubeny--Rangamani--Takayanagi entropy-area rule.
Its declared assumptions are: an applicable AdS/CFT dictionary, the relevant
asymptotically AdS boundary class, and a state class where the RT/HRT
semiclassical entropy-area prescription is valid.
\end{definition}

\begin{definition}[Finite diagnostic carrier]
At refinement level \(r\in\{2,4\}\), let the finite carrier be
\[
  E_r=\mathbb{R}^{r+2}
\]
with audit energy \(C=I\).  The first \(r\) coordinates are boundary
entanglement-ledger coordinates \(s_1,\ldots,s_r\).  The last two coordinates
are \(g\), an independent bulk geometry-amplitude coordinate, and \(m\), an
independent P3 route-mismatch coordinate.

The native recognition lens \(L_r\) contains only the \(s_i\) rows.  The
RT-like dissolving area rows are \(a_i=s_i/\sqrt r\).  The full E018 diagnostic
dissolving family is
\[
  D_r=\{a_1,\ldots,a_r,\; g+0.25\,\bar s,\; m+0.30\,g+0.10\,\bar s\},
  \qquad
  \bar s=r^{-1/2}\sum_i s_i .
\]
\end{definition}

\begin{proposition}[Finite-carrier Mode A obstruction]
On the declared carrier,
\[
  \Xi_I(D_{\mathrm{area}}\mid L_r)=0
\]
for \(r=2,4\), while the full joint target has
\[
  \tr\,\Xi_I(D_r\mid L_r)=2.09,
  \qquad
  \frac{\tr\,\Xi_I(D_r\mid L_r)}{\tr K_{DD}}=0.6608695652173913
\]
at both refinement levels.
\end{proposition}

\begin{proof}
The RT-like area rows lie in the span of the native entanglement-ledger rows,
so exact adequacy holds for the area sector by factorization through \(L_r\).
The bulk-amplitude coordinate \(g\) and the P3 coordinate \(m\) are orthogonal
to the native span.  The Schur complement therefore removes only the
entanglement-ledger components of the last two dissolving rows and leaves the
positive residual carried by \(g\) and \(m\).  Direct computation in
\texttt{simulate\_step1\_holographic\_xi.py} gives the trace and normalized
values displayed above.
\end{proof}

\begin{remark}[Overread control]
If \(g\) and \(m\) are added to the native lens, the full diagnostic residual
vanishes.  Step 1 rejects that variant because it adds the target answer rows
to the native recognition source.  The vanishing residual in that overread
control is therefore not a closure witness.
\end{remark}

\begin{remark}[Recognition verdict]
The leading import supplies a partial E2 recognition source for the
AdS/RT-like area-shadow sector.  It does not land the canonical E018 root:
the unrestricted amplitude-over-geometry object and the general P3
route-mismatch remain outside the audited import.
\end{remark}

\end{document}
